\section{Classical Radar Emitter Identification}\label{sec:related:classical}

Operational \gls{elint} systems recover emitter identity from measured \glspl{pdw} by combining pulse sorting with library lookup \parencite{wiley2006elint,skolnik2008radar}.
After detection, pulses are grouped into trains and each train is matched against stored templates of carrier frequency, pulse width, \gls{pri} family, and scan behavior.
The grouping step is deinterleaving; the matching step is identification.
Both historically depend on hand-specified features and on regularity that agile radars no longer guarantee.

The oldest deinterleavers exploit repetition in \gls{toa}.
The cumulative difference histogram (CDIF) accumulates \gls{toa} differences at increasing pulse lags and treats histogram peaks as candidate \glspl{pri}, which a subsequent search then validates as pulse sequences \parencite{mardia1989deinterleaving}.
The sequential difference histogram (SDIF) keeps only the current difference level, with a derived detection threshold, which avoids accumulating every lag while remaining a \gls{toa}-only method \parencite{milojevic1992sdif}.
These algorithms are reliable when each emitter keeps an approximately constant \gls{pri} and when missing pulses and cross-emitter differences do not create false peaks.
They degrade under stagger, jitter, and dense mixing, which is the setting of the recordings in \cref{sec:method:formulation}.

When timing alone is unstable, sorting moves into a joint space of \gls{rf}, pulse width, and \gls{aoa}.
Centroid methods such as \(k\)-means and density clustering such as \gls{dbscan} then replace histogram peak finding \parencite{ester1996density,qu2026radardeinterleavingreview}.
Density clustering has the operational advantage that the number of emitters need not be specified in advance.
Its accuracy still depends on the metric, the scaling of each coordinate, and a neighbourhood radius that is hard to set when cluster densities differ.
Offline direction clustering of overlapping pulses is one concrete instance of this approach \parencite{bedoire2022offlineclustering}.
The method chapter uses \gls{dbscan} only after a learned embedding; applying it to raw window features is retained as a baseline (\cref{sec:method:baselines}).

Once trains are formed, identification is typically template matching against a mode or equipment library \parencite{wiley2006elint}.
Statistical fingerprints---histograms of \gls{rf} and \gls{pri}, dwell statistics, scan-period estimates---summarize each train.
The match assumes that the library contains the emitter and that deinterleaving errors have not corrupted the fingerprint.
Novel emitters and modes that are absent from the catalogue fall outside this procedure by construction.

Classical pipelines treat deinterleaving and mode analysis as successive stages, each with its own hand-engineered features.
They do not learn a representation in which both partitions are visible, and they have no mechanism for emitters that appear only at inference.
Those two requirements---recording-local emitter grouping and a mode grouping that is not a closed-catalogue softmax---motivate the learned encoder of \cref{ch:method}.
